\clearpage
\scanpage{83}
\reversemarginpar
\margin{\figure{\tikz[x=1em,y=1em]{\draw[line width=.22pt] (0,0) -- (1.48,-.04); \draw[black!48,line width=.2pt] (1.48,-.04) -- (2.02,-.05); \fill (0,0) ellipse (.13 and .055); \fill (1.48,-.03) ellipse (.045 and .033); \draw (2,-.01) -- (2.03,-.09); \draw[black!50,line width=.26pt] (.40,.64) .. controls (.51,.61) and (.50,.58) .. (.43,.53); \draw[black!22,line width=.18pt] (.44,.51) -- (.40,.04);}}}
\normalmarginpar
\mspage{B-1}

\begin{center}
ELLIPTIC CURVES :\quad BIBLIOGRAPHY
\end{center}

1)\quad \uline{Elliptic and modular functions}

C.JORDAN - Course in Analysis, Gauthi\makebox[0pt][l]{\deleted{\unsure{o|e}}}er-Villars

A.HURWITZ and R.COURANT - Vorlesungen über allgemeine Funktionentheorie und elliptische Funktionen, Springer (4th edition - 1964)

R.FRICKE - Die elliptischen Funktionen und ihre Anwendungen, Teubner, Leipzig, 1922

2)\quad \uline{Complex multiplication} (\uline{analytic methods})

H.WEBER - Algebra, volume III, Strasbourg, 1908

R.FUETER - Vorlesungen über die singulären Moduln und die komplexe Multiplik.tion der elliptischen Funktionen, I, 1924, II, 1927, Teubner.

H.HASSE - Neue Begründung der komplexen Multiplikation, I, Crelle's Journal, 157, pages 115-139 (1927); II, \uline{in the same place}, 165, pages 64-88 (1931)

Seminar on complex multiplication, Inst.Adv.Study, Princeton, 1957-1958

M.DEURING - Die Klassenkörper der komplexen Multiplikation, Encyclopaedia of Mathematical Sciences, I 2, 23, Teubner, 1958

3)\quad \uline{Algebraic function fields of one variable}

C.CHEVALLEY - Introduction to the theory of algebraic functions of one variable, Math.Surveys, VI, New-York, 1951

M.EICHLER - Einführung in die Theorie der algebraischen Zahlen und Funktionen, Birkhäuser Verlag, 1963
% archive sheet number: 82
